% !TeX root = ../../Book1.tex
% 纲要标题：### 20.4 SU(2) 的最高权分类与球面线丛实现
% 纲要位置：几何学计划纲要.md，第 912 行。
% 正文状态：初稿；尚未正式审核。

% 纲要内容开始
% 
% 以 SU(2) 的权分解和球面线丛说明后续最高权及全纯实现的问题；明确哪些结论已由本章证明，哪些需要第四卷复几何。
% 
% ---
% 
% 纲要内容结束

\section{SU(2) 的最高权分类与球面线丛实现}
\label{g1:sec:2004}

Peter--Weyl 定理告诉我们所有不可约表示的矩阵系数合在一起足够丰富，却没有逐个列出这些表示。上一节又把表示与齐性丛截面联系起来，但一个光滑截面空间通常无限维。现在要分别解决两个更具体的问题：怎样用权辨认不可约表示，以及怎样在一个截面空间中找出自然的有限维表示。我们在 \(\operatorname{SU}(2)\) 上把这两步都算出来，再说明一般紧群所需的新结构。

\ProseParagraphBreak

固定
\[
	T=\{t(\theta)=\operatorname{diag}(e^{\mathrm{i}\theta},e^{-\mathrm{i}\theta}):\theta\in\mathbb R\},
	\qquad k=\begin{pmatrix}\mathrm{i}&0\\0&-\mathrm{i}\end{pmatrix}.
\]
环面的周期为 \(2\pi\)，所以它的特征标恰为 \(t(\theta)\mapsto e^{\mathrm{i}m\theta}\)、\(m\in\mathbb Z\)。若 \(V\) 是有限维复酉表示，将 \(\mathrm{d}\pi(k)\) 对角化，就得到
\[
	V=\bigoplus_{m\in\mathbb Z}V[m],\qquad
	V[m]=\{v:\pi(t(\theta))v=e^{\mathrm{i}m\theta}v\text{ 对所有 }\theta\}.
\]
这里只出现有限多个非零空间。整数性也可直接核验：斜 Hermitian 算子的特征值写成 \(\mathrm{i}a\)，而 \(t(2\pi)=I\) 迫使 \(e^{2\pi\mathrm{i}a}=1\)。本节将整数 \(m\) 简称为相应的权；它代表 Lie 代数上的实泛函 \(\theta k\mapsto m\theta\)。第十九章的正根为 \(\alpha(\theta k)=2\theta\)，所以沿根方向改变权时，整数标签改变二。

\ProseParagraphBreak

先构造候选表示。对整数 \(n\ge0\)，记
\[
	V_n=\{p(z_0,z_1):p\text{ 是总次数 }n\text{ 的复齐次多项式}\},
	\qquad (\pi_n(g)p)(z)=p(g^{-1}z).
\]
\symidx{\(V_n,\pi_n\)：\(\operatorname{SU}(2)\) 的 \(n\) 次齐次多项式表示}常数构成 \(V_0\)。逆矩阵保证这是左表示：依次代入 \(g^{-1}\)、\(h^{-1}\) 得到 \((gh)^{-1}\)，故 \(\pi_n(g)\pi_n(h)=\pi_n(gh)\)。系数光滑依赖 \(g\)，所以它是光滑表示。其维数为 \(n+1\)，由平均内积可以酉化。

\begin{theorem}[特殊酉二阶群的不可约表示]\label{g1:thm:su2-polynomial-classification}
	取 \(T=\{t(\theta)=\operatorname{diag}(e^{\mathrm{i}\theta},e^{-\mathrm{i}\theta}):\theta\in\mathbb R\}\subseteq\operatorname{SU}(2)\)。在表示 \(\pi\) 中，约定非零向量 \(v\) 满足 \(\pi(t(\theta))v=e^{\mathrm{i}m\theta}v\)（所有 \(\theta\in\mathbb R\)）时称其为权 \(m\) 向量，其中 \(m\in\mathbb Z\)。对每个整数 \(n\ge0\)，二元 \(n\) 次复齐次多项式空间 \(V_n\) 配备作用 \((\pi_n(g)p)(z)=p(g^{-1}z)\)，是 \(\operatorname{SU}(2)\) 的不可约复表示，权为
	\[
		n,n-2,\ldots,-n,
	\]
	每个权空间均为一维。每个有限维不可约复光滑表示都与唯一一个 \(V_n\) 等价。\(n\) 是该表示的\term[zuigaoquan]{最高权}[highest weight]。
\end{theorem}
\begin{proof}
	先在候选表示中实际计算升降算子。\(\mathfrak{su}(2)\) 的复化为 \(\mathfrak{sl}(2,\mathbb C)\)，取
	\[
		e=\begin{pmatrix}0&1\\0&0\end{pmatrix},\quad
		f=\begin{pmatrix}0&0\\1&0\end{pmatrix},\quad
		h=\begin{pmatrix}1&0\\0&-1\end{pmatrix}.
	\]
	它们满足 \([h,e]=2e\)、\([h,f]=-2f\)、\([e,f]=h\)，且 \(k=\mathrm{i}h\)。对 \(p(g^{-1}z)\) 微分，得到
	\[
		E=-z_1\frac{\partial}{\partial z_0},\qquad
		F=-z_0\frac{\partial}{\partial z_1},\qquad
		L=-z_0\frac{\partial}{\partial z_0}+z_1\frac{\partial}{\partial z_1}.
	\]
	因此 \(z_0^jz_1^{n-j}\) 的权为 \(n-2j\)，且 \(F\) 将它送到 \(-(n-j)z_0^{j+1}z_1^{n-j-1}\)，\(E\) 则向相反方向移动，系数为 \(-j\)。两端以外的这些系数都非零。

	设 \(W\subseteq V_n\) 是非零群不变子空间。微分后，它在 \(E,F,L\) 下也不变。\(L\) 有互不相同的特征值，每个特征空间投影都可写成 \(L\) 的插值多项式，所以 \(W\) 包含某个单项式。反复使用 \(E,F\) 就得到全部单项式，故 \(W=V_n\)。这证明不可约性。

	反过来，设 \(\pi\) 是任一不可约表示，先用\cref{g1:prop:compact-unitarization-complete-reducibility}选不变 Hermitian 内积。把微分复线性延拓，记 \(E=\mathrm{d}\pi(e)\)、\(F=\mathrm{d}\pi(f)\)、\(L=\mathrm{d}\pi(h)\)。实 Lie 代数中的
	\[
		x=e-f,\qquad y=\mathrm{i}(e+f),\qquad k=\mathrm{i}h
	\]
	在酉表示中微分为斜 Hermitian 算子。由 \(e=(x-\mathrm{i}y)/2\)、\(f=(-x-\mathrm{i}y)/2\)、\(h=-\mathrm{i}k\)，可得 \(E^*=F\)、\(L^*=L\)。三个交换关系也由微分保持括号而成立。因此\cref{g1:lem:sl2-weight-chain}的条件全部满足，空间分成有限条升降链。

	每条链都在 \(E,F,L\) 下不变，因而在实 Lie 代数下不变，再由指数自然性在每个 \(\pi(\exp X)\) 下不变。\(\operatorname{SU}(2)\) 连通，\cref{g1:prop:exponentials-generate-identity-component}说明指数生成整个群，所以每条链也是群不变子空间。不可约性迫使只有一条链。于是存在非负整数 \(n\) 及非零向量 \(v\)，使
	\[
		Ev=0,\quad Lv=nv,\quad
		v,Fv,\ldots,F^nv
	\]
	为一组基，且 \(E(F^jv)=j(n-j+1)F^{j-1}v\)。这正是候选表示中从 \(z_1^n\) 开始的链所满足的公式。把 \(F^jv\) 送到 \(F^jz_1^n\)，得到与三个算子交换的线性同构。它因而与所有指数的表示交换，再与整个群交换，给出 \(\pi\cong\pi_n\)。不同 \(n\) 的维数不同，唯一性随之成立。
\end{proof}

最高权来自升降链的终止。选定正根以后，升算子不断把权加二；有限维性保证这个过程终止，链顶被升算子杀掉。交换关系又强迫链顶是非负整数，并准确决定向下还能走多少步。一般紧群有多个根方向，这才出现如何让不同升降方向相容的新问题。

\begin{corollary}[特征标与旋转群表示]\label{g1:cor:su2-character-so3-descent}
	设 \(n\ge0\) 为整数，\(V_n\) 为二元 \(n\) 次复齐次多项式空间，取 \(\operatorname{SU}(2)\) 作用 \((\pi_n(g)p)(z)=p(g^{-1}z)\)，并令 \(t(\theta)=\operatorname{diag}(e^{\mathrm{i}\theta},e^{-\mathrm{i}\theta})\)。则 \(V_n\) 的特征标在该对角环面上为
	\[
		\chi_n(t(\theta))=\sum_{j=0}^n e^{\mathrm{i}(n-2j)\theta}
		=\frac{\sin((n+1)\theta)}{\sin\theta}\qquad(\sin\theta\ne0).
	\]
	在 \(\theta=0\) 时其值为 \(n+1\)，在 \(\theta=\pi\) 时为 \((-1)^n(n+1)\)。该表示下降到 \(\operatorname{SO}(3)\cong\operatorname{SU}(2)/\{\pm I\}\) 当且仅当 \(n\) 为偶数。因此 \(\operatorname{SO}(3)\) 的不可约复光滑表示恰由非负整数 \(\ell\) 标记，维数为 \(2\ell+1\)。
\end{corollary}
\begin{proof}
	迹等于各一维权空间上的特征值之和，有限等比级数求和给出正弦商；分母为零时直接使用有限和，不把商式当作定义。中心元 \(-I\) 作用为 \(p(z)\mapsto p(-z)=(-1)^np(z)\)。由\cref{g1:prop:quaternion-low-rank}的二重覆盖，下降恰要求这个核作用平凡，即 \(n\) 为偶数。下降后的表示光滑，因为覆盖有光滑局部逆。反过来，将旋转群的不可约表示沿该满射拉回，仍不可约，且核作用平凡；分类定理于是给出唯一的 \(n=2\ell\)。
\end{proof}

% !TeX root = ../Book1.tex
\begin{figure}[htbp]
	\centering
	\begin{tikzpicture}[font=\small,x=0.95cm,y=1cm]
		\node[left] at (-0.3,1) {\(V_4\)};
		\foreach \x/\m in {0/4,2/2,4/0,6/-2,8/-4}{
			\fill[GeometryBlue] (\x,1) circle (2pt);
			\node[above] at (\x,1.13) {\(\m\)};
		}
		\foreach \x in {0,2,4,6}{
			\draw[->,GeometryBlue,thick] (\x+0.15,1)--(\x+1.85,1) node[midway,below] {\(F\)};
		}
		\node[right] at (8.35,1) {下降到 \(\operatorname{SO}(3)\)};
		\node[left] at (-0.3,-0.7) {\(V_3\)};
		\foreach \x/\m in {1/3,3/1,5/-1,7/-3}{
			\fill[GeometryBlue] (\x,-0.7) circle (2pt);
			\node[above] at (\x,-0.57) {\(\m\)};
		}
		\foreach \x in {1,3,5}{
			\draw[->,GeometryBlue,thick] (\x+0.15,-0.7)--(\x+1.85,-0.7) node[midway,below] {\(F\)};
		}
		\node[right] at (8.35,-0.7) {\(-I\) 作用为 \(-1\)};
	\end{tikzpicture}
	\caption[特殊酉二阶群的权链]{\(V_4\) 与 \(V_3\) 的每个权空间都是一维。降算子 \(F\) 每次把权减二，升算子 \(E\) 沿相反方向作用；图中箭头对应非零作用，\(F\) 在最低权处为零，\(E\) 在最高权处为零。偶数权链使中心元 \(-I\) 作用平凡，奇数权链则不满足下降条件。}
	\label{g1:fig:su2-weight-chains}
\end{figure}


\cref{g1:fig:su2-weight-chains}显示了权的奇偶性如何记录整体群。相同 Lie 代数上的升降公式不能单独保证表示下降；必须检查周期格，也即这里的中心元 \(-I\)。这正是第十九章根数据与整体群区别在表示论中的一个可计算后果。

现在回到球面上的线丛。沿用\cref{g1:prop:su2-hopf-associated-lines}的记号，\(E_m=\operatorname{SU}(2)\times_{\chi_m}\mathbb C\to\mathbb{CP}^1\)，其中 \(\chi_m(t(\theta))=e^{\mathrm{i}m\theta}\)。\(E_1\) 是重言线丛，纤维为点 \(\ell\in\mathbb{CP}^1\) 所代表的直线本身；\(E_{-n}\) 是它的 \(n\) 次对偶张量幂。为了说明多项式怎样变成截面，先把每条直线上的齐次性解释为该纤维中的一个元素。

\begin{proposition}[球面线丛中的多项式表示]\label{g1:prop:polynomial-sections-sphere-line}
	设 \(n\ge0\) 为整数，\(V_n\) 为二元 \(n\) 次复齐次多项式空间，取 \(\operatorname{SU}(2)\) 作用 \((\pi_n(g)p)(z)=p(g^{-1}z)\)。令 \(T=\{t(\theta)=\operatorname{diag}(e^{\mathrm{i}\theta},e^{-\mathrm{i}\theta}):\theta\in\mathbb R\}\subseteq\operatorname{SU}(2)\)；对 \(j\in\mathbb Z\)，置 \(\chi_j(t(\theta))=e^{\mathrm{i}j\theta}\)、\(E_j=\operatorname{SU}(2)\times_{\chi_j}\mathbb C\)，采用 \([gt,c]=[g,\chi_j(t)c]\) 的商约定。每个 \(p\in V_n\) 确定 \(E_{-n}\) 的光滑截面，等变函数为
	\[
		f_p(g)=p(ge_0),\qquad e_0=(1,0)^{\mathsf T}.
	\]
	这一映射 \(V_n\to\Gamma^\infty(E_{-n})\) 单射且 \(\operatorname{SU}(2)\) 等变。在 \(U_0=\{z_0\ne0\}\)、\(U_1=\{z_1\ne0\}\) 的坐标 \(z=z_1/z_0\)、\(w=z_0/z_1\) 中，相对于 \((1,z)\)、\((w,1)\) 的对偶张量标架，截面分量分别为
	\[
		a_0(z)=p(1,z),\qquad a_1(w)=p(w,1),\qquad
		a_0(z)=z^n a_1(1/z).
	\]
	它的像恰为这两张图中分量都为多项式的截面。若将 \(n\) 换成负整数，则满足相同多项式过渡条件的截面只有零截面。
\end{proposition}
\begin{proof}
	齐次性使 \(p\) 在一条直线 \(\ell\) 上的限制成为 \((\ell^*)^{\otimes n}\) 的一个元素：若选 \(v\ne0\) 作基，只需规定它在 \(v^{\otimes n}\) 上取值 \(p(v)\)；把基改成 \(cv\) 时，输入和取值都乘 \(c^n\)，故与选择无关。把这个构造放在各条直线上，就得到所述截面。在两张图中写出的分量都是多项式，因此光滑。

	用单位标架 \(ge_0\) 表示同一个截面，所得函数正是 \(f_p\)。右乘 \(t(\theta)\) 使第一列乘以 \(e^{\mathrm{i}\theta}\)，故
	\[
		f_p(gt(\theta))=e^{\mathrm{i}n\theta}f_p(g)
		=\chi_{-n}(t(\theta))^{-1}f_p(g),
	\]
	符合\cref{g1:prop:equivariant-functions-sections}的条件。若截面为零，则 \(p\) 在所有单位向量上为零；齐次性使它在全部非零向量上为零，也在零点为零，故 \(p=0\)。对 \(a\in\operatorname{SU}(2)\)，
	\[
		f_{\pi_n(a)p}(g)=p(a^{-1}ge_0)=f_p(a^{-1}g),
	\]
	正是截面上的左作用，所以映射等变。

	在交集上，\((w,1)=w(1,z)\)、\(w=1/z\)，故对偶的 \(n\) 次张量幂给出 \(a_0(z)=z^n a_1(1/z)\)。若 \(a_1(w)=\sum_{j=0}^N b_jw^j\)，右侧为 \(\sum_j b_jz^{n-j}\)。它与多项式 \(a_0\) 在 \(\mathbb C^*\) 上相等，说明所有负次幂的系数都为零。对 \(n\ge0\)，只有 \(0\le j\le n\) 可以出现，而
	\[
		p(z_0,z_1)=\sum_{j=0}^n b_j z_0^jz_1^{n-j}
	\]
	恰好产生这两个分量。对 \(n<0\)，所有可能的次数 \(n-j\) 都为负，只能全部系数为零。这里的比较只是有限 Laurent 多项式系数的唯一性，无须对一般解析函数作展开。
\end{proof}

\(E_{-n}\) 的光滑截面远不止这些多项式截面。实际上，每一种奇偶性合适、足够大的不可约表示都会出现。以下结果精确说明两种空间的差别。

\begin{proposition}[球面线丛的光滑表示重数]\label{g1:prop:sphere-line-section-multiplicity}
	设 \(m\in\mathbb Z\)，\(V_n\) 为二元 \(n\) 次复齐次多项式表示，\(n\ge0\) 为整数，作用为 \((\pi_n(g)p)(z)=p(g^{-1}z)\)。令 \(T=\{t(\theta)=\operatorname{diag}(e^{\mathrm{i}\theta},e^{-\mathrm{i}\theta}):\theta\in\mathbb R\}\subseteq\operatorname{SU}(2)\)，\(\chi_m(t(\theta))=e^{\mathrm{i}m\theta}\)，\(E_m=\operatorname{SU}(2)\times_{\chi_m}\mathbb C\)，采用 \([gt,c]=[g,\chi_m(t)c]\) 的商约定。截面上取左作用 \((U_a\sigma)(x)=a\cdot\sigma(a^{-1}x)\)，其中 \(a\cdot[g,c]=[ag,c]\)。则
	\[
		\dim\operatorname{Hom}_{\operatorname{SU}(2)}
		\bigl(V_n,\Gamma^\infty(E_m)\bigr)
		=\begin{cases}
			1,& n\ge|m|\text{ 且 }n\equiv m\pmod2,\\
			0,&\text{其他情形}.
		\end{cases}
	\]
	把光滑截面按不变 Hermitian 度量和归一化商体积作 \(L^2\) 完备化，所得 Hilbert 空间 \(\mathcal H_m\) 在左群作用下分解为
	\[
		\mathcal H_m\cong
		\widehat{\bigoplus}_{\substack{n\ge|m|\\n\equiv m\ (2)}}V_n,
	\]
	每个表示恰出现一次。这是 Hilbert 空间中的收敛陈述，不表示任意光滑截面的任意排序级数都一致或逐阶可微收敛。
\end{proposition}
\begin{proof}
	\cref{g1:thm:smooth-frobenius-reciprocity}把所求同态空间识别为 \(\operatorname{Hom}_T(V_n,\mathbb C_m)\)，其中 \(T\) 在 \(\mathbb C_m\) 上按 \(\chi_m\) 作用。一个这样的线性映射在权 \(q\ne m\) 的空间上必为零：取使 \(e^{\mathrm{i}q\theta}\ne e^{\mathrm{i}m\theta}\) 的 \(\theta\)，等变性就迫使该空间的像为零。在权 \(m\) 空间上则可任取线性泛函。\(V_n\) 的权是 \(n,n-2,\ldots,-n\)，各重数一，故得到第一式。

	再说明为何这份重数表穷尽整个 \(L^2\) 空间。令 \(\mathrm{d}t\) 为 \(T\) 的归一化积分，在 \(L^2(\operatorname{SU}(2))\) 上定义
	\[
		(P_m f)(g)=\int_T\chi_m(t)f(gt)\,\mathrm{d}t.
	\]
	积分可先对连续函数定义，再由 \(\|P_mf\|_2\le\|f\|_2\) 延拓；这个估计来自 Cauchy--Schwarz 不等式与右不变换元。对连续函数换元 \(u=st\) 得 \((P_mf)(gs)=\chi_m(s)^{-1}(P_mf)(g)\)，而已经满足该条件的函数被 \(P_m\) 固定。因此 \(P_m^2=P_m\)。酉性、取逆不变性及 \(\overline{\chi_m(t)}=\chi_m(t^{-1})\) 给出 \(P_m^*=P_m\)，故它是正交投影，并与所有左平移交换。

	\cref{g1:thm:peter-weyl-hilbert}给出矩阵系数有限和在群上 \(L^2\) 稠密。\(P_m\) 把每个有限维系数空间仍送回该空间，所得函数光滑且满足上述等变性，所以其像正是 \(\mathcal H_m\) 中的一个稠密子空间。等变函数的模沿右 \(T\) 陪集不变，群积分下推到商上的积分，因而该识别保持范数。不同不可约类型正交，由第一部分给出的重数，每个允许的 \(V_n\) 恰留下一份，其他类型消失，得到 Hilbert 直和。这里也可注意 \(V_n\) 的对偶仍不可约，权集合取负后不变，因此按分类仍等价于 \(V_n\)；使用左正则表示或其对偶的标记不会改变此处的重数表。
\end{proof}

\begin{example}\label{g1:exm:sphere-sections-smooth-polynomial}
	\(m=0\) 时，线丛平凡，光滑截面就是球面上的光滑函数。上式给出维数依次为 \(1,3,5,\ldots\) 的不可约部分。常数构成 \(V_0\)；在单位球面 \(S^2\subseteq\mathbb R^3\) 上，三个坐标函数的复线性张成空间构成三维的 \(V_2\)。后一断言可由第十九章 \(\operatorname{SU}(2)\) 在虚四元数上的共轭作用核验：其复化是 \(\mathfrak{sl}(2,\mathbb C)\) 的三维升降链，因而对应 \(n=2\)。

	对 \(m=1\)，重言线丛虽然没有处处非零的光滑截面，却有一个二维不变截面空间。给定 \(a\in\mathbb C^2\)，令 \(\sigma_a(\ell)\) 为 \(a\) 在直线 \(\ell\) 上的正交投影。它光滑依赖 \(\ell\)，且酉变换保持正交投影，故 \(a\mapsto\sigma_a\) 等变。取 \(\ell=\mathbb Ca\) 可知它单射；非零 \(a\) 对应的截面恰在 \(\ell=a^\perp\) 处为零。这与丛的非平凡性相容。与此同时，前一命题在 \(n=-1\) 的结论说明 \(E_1\) 没有非零的两图多项式截面。光滑条件与多项式条件已经选出了不同的表示空间。
\end{example}

由此可以看清一般理论下一步要增加什么。权是极大环面的特征标，根描述不同权空间之间可由哪些 Lie 代数方向连接；最高权还依赖正根的选择。\(\operatorname{SU}(2)\) 只有一个正根方向，有限升降链已经给出完整分类。对于一般紧连通 Lie 群，需要研究多个根方向的相容性，并同时保留特征标格对整体群的限制；一般最高权分类在这一层面组织不可约表示。

\ProseParagraphBreak

几何方面，多项式截面给出了 \(V_n\) 的显式实现，但选取更一般的截面类别还需要底空间及线丛的复结构。Borel--Weil 理论\footnote{Armand Borel（1923--2003），瑞士数学家，研究代数拓扑与线性代数群。André Weil（1906--1998），法国数学家，对代数几何与数论有重要贡献；其姓氏 Weil 与本章 Weyl 是不同姓名。}研究在旗流形的全纯线丛截面中实现表示。第四卷将建立全纯坐标、全纯丛与相应的解析工具，这些概念使“哪些截面应当保留”成为内在的几何条件。在已经算出的秩一模型中，整数权决定线丛的过渡函数，多项式条件则选出了相应的有限维表示。
% 本节沿用纲要的连续叙述安排，不另造小节标题。
